\chapter{Incomplete Factorization and Domain Decomposition}\label{v0ch:18}
\begin{chaptersummary}
Sequential triangular solves and sweeps do not scale directly across many processors. Domain decomposition partitions the unknowns into subdomains, combining local solves with overlapping data exchange to construct parallel preconditioners. Once local coupling is resolved, global low-frequency error still requires a coarse space—leading directly to multigrid.
\end{chaptersummary}
\learninggoals{Write additive Schwarz; distinguish overlap, RAS, and symmetry; explain the need for a coarse space.}
\paragraph{Core Question}How can approximate inverse action be parallelized without losing essential global information?

\section{Combining Local Solves}
Let $R_i$ restrict a global vector to a subdomain, with $A_i=R_iAR_i^T$. Additive Schwarz applies
\[
 B=\sum_i R_i^TA_i^{-1}R_i,\qquad \bm z=B\bm r.
\]
Subdomains may overlap, and local problems can be solved with exact LU, ILU, or other approximations. When $A$ is SPD, the subdomains cover the domain, and each local inverse action is exact or SPD, the additive Schwarz operator $B$ is SPD; multiplicative Schwarz updates residuals sequentially in analogy with subdomain Gauss--Seidel.
Restricted additive Schwarz (RAS) injects local corrections only into nonoverlapping interior cores, reducing communication redundancy at the cost of a generally nonsymmetric preconditioner action—a structural distinction relevant when pairing RAS with standard PCG.
\section{Coarse Spaces and Block Systems}
Single-level subdomain preconditioners propagate smooth error slowly across many partitions. Adding a coarse space $Z$ via
\[
 B_0=Z(Z^TAZ)^{-1}Z^T,
\]
addresses global low-frequency modes alongside the local corrections. Without a coarse level, iteration counts typically grow with the number of subdomains, with precise bounds governed by overlap width, subdomain size, and the PDE.
For saddle-point systems, block elimination produces the Schur complement
$S=A_{22}-A_{21}A_{11}^{-1}A_{12}$, enabling approximate block-diagonal or block-triangular preconditioners that respect distinct physical scales rather than treating all fields with a single generic ILU.
\conceptbridge{Why can accurate local solutions still yield slow global convergence?}{Exact local solves remove internal subdomain error, yet slowly varying modes spanning multiple subdomains can survive every local correction. A coarse space provides the global coupling mechanism that increasing local solve accuracy cannot replace.}{Accurate local solve $\longrightarrow$ remaining interdomain modes $\longrightarrow$ global coarse correction.}
\section{A Minimal Experiment}
Partitioning one-dimensional Poisson into four subdomains compares block Jacobi, one-layer-overlap additive Schwarz, and a two-level variant with a subdomain-constant coarse space. Scaling the subdomain count at fixed total degrees of freedom while recording local setup, halo exchange volume, and iteration counts isolates the impact of coarse correction, with multiprocess timings quantifying distributed scaling.
\costnote{Wider overlap raises local storage and communication overhead while often lowering iteration counts. Local ILU confines triangular-solve dependencies within subdomains, shifting potential scalability bottlenecks to the coarse-grid solve and global reductions.}
\aiexercise{For eight-point Poisson, write two overlapping restrictions $R_i$ and $B$. Check symmetry, then replace injection by RAS and check again.}{Derive the Schur complement of a pressure--velocity system and explain why block-triangular preconditioning often accompanies GMRES.}
\studyroute{Schwarz condition-number bounds, BDDC, and FETI-DP are omitted. See \cite{saad,petsc}.}{additive Schwarz; restricted additive Schwarz; domain decomposition; coarse space; Schur complement}
